% ═══════════════════════════════════════════════════════════════════════════════
% PAPER STRUCTURE TEMPLATE (MACRO TRACK) — Auto-filled by Stage 5
%
% Model-based general-equilibrium paper (Kaplan-Moll-Violante / Achdou et al.
% style). The LaTeX preamble is loaded separately from example.tex.
%
% RULES:
% 1. Do NOT reorder sections; do NOT delete sections
% 2. Every {{PLACEHOLDER}} must be replaced with actual content
% 3. Numbers MUST come from data/clean/main_results.csv and data/model/*.json
% 4. No causal-inference vocabulary: there is no "treatment", "identification
%    strategy" or "estimate" here — the model is solved and calibrated
% ═══════════════════════════════════════════════════════════════════════════════

\title{ {{PAPER_TITLE}} }
\author{ {{AUTHOR_LINE}} }
\date{\today}

\begin{document}
\maketitle

\begin{abstract}
\noindent
{{ABSTRACT_QUESTION}}
% ONE sentence: the macro question and why heterogeneity / GE matters for it.
{{ABSTRACT_MODEL}}
% ONE sentence: the model (class, key friction) and how it is disciplined by data.
{{ABSTRACT_RESULTS}}
% 2 sentences with exact numbers: main quantitative result and comparison with
% the RANK/TANK benchmark (e.g. impact response, direct-effect share).
{{ABSTRACT_IMPLICATIONS}}
\end{abstract}

\textbf{JEL Codes:} {{JEL_CODES}} \\
\textbf{Keywords:} {{KEYWORDS}}

\newpage

\section{Introduction}
\label{sec:introduction}
% P1 question & motivation; P2 the mechanism in one sentence; P3 the model and
% its discipline (targeted moments); P4 main results with numbers; P5 why the
% answer differs from RANK/TANK; P6 contribution to literature; P7 road map.
{{INTRO}}

\section{Related Literature}
\label{sec:literature}
% (i) HA macro / HANK; (ii) the topic-specific literature; (iii) methods
% (continuous time, sequence space). Cite only papers you are certain exist.
{{LITERATURE}}

\section{Model}
\label{sec:model}
\subsection{Environment}
{{MODEL_ENVIRONMENT}}
% Agents, preferences, income process (Poisson types, generator Lambda), assets
% and borrowing limit, technology, government budget, monetary/fiscal rule.
\subsection{Household problem}
{{MODEL_HOUSEHOLD}}
% HJB equation and KFE in continuous time (Achdou et al. 2022), state
% constraint at the borrowing limit, first-order condition.
\subsection{Equilibrium}
{{MODEL_EQUILIBRIUM}}
% Formal definition (prices, policies, distribution, market clearing).
% Remark: which market is redundant by Walras' law and why (aggregate budget).

\section{Analytical Results}
\label{sec:analytical}
% Propositions with proofs (in the appendix): e.g. r < rho with a borrowing
% limit; TANK multiplier; decomposition of the response into direct and
% indirect effects. If Stage 5.5 (Lean) ran, state which results are
% machine-checked and which remain informal, exactly as lean/ reports.
{{ANALYTICAL_RESULTS}}

\section{Calibration}
\label{sec:calibration}
% Externally set vs internally calibrated parameters (Table tab:calibration);
% targeted and untargeted moments (Table tab:moments). Explain which moment
% pins which parameter and why that moment matters for the mechanism.
{{CALIBRATION}}

\section{Results}
\label{sec:results}
\subsection{Stationary equilibrium}
{{RESULTS_STEADY_STATE}}
\subsection{Main experiment}
{{RESULTS_EXPERIMENT}}
% Figures fig_irf_output / fig_irf_decomposition or fig_counterfactuals;
% Table tab:experiment. Compare HANK with TANK and RANK and explain the gap.

\section{Sensitivity}
\label{sec:sensitivity}
% How the main number moves with the parameters that drive it; grid and
% horizon robustness; every SOFT test that failed is discussed here.
{{SENSITIVITY}}

\section{Conclusion}
\label{sec:conclusion}
{{CONCLUSION}}
% Include an honest limitations paragraph (what the model abstracts from).

\appendix
\section{Computational Appendix}
\label{app:computation}
% Algorithm (implicit upwind finite differences; KFE as the transpose of the
% generator; sequence-space Newton), grids and tolerances, and Table
% tab:eqtests with every equilibrium test (Walras' law, market clearing, ...).
{{COMPUTATIONAL_APPENDIX}}

\section{Proofs}
\label{app:proofs}
{{PROOFS}}

% thebibliography goes here (see LaTeX rules)
\end{document}
